% !TeX root = ../drafts/08-end-to-end-protocol.tex
\section{End-to-end protocol}\label{sec:e2e}

\subsection{Bytecode encoding}\label{sec:e2e-bc}

The bytecode (of $\nprog=2^{\kbc}$ instructions) is encoded as one $\K$-valued multilinear $P$ in $\kbc+4$ variables:
\begin{enumerate}[itemsep=1pt,topsep=3pt]
\item lay instruction $z$ out as $16$ slots of $\K$: the opcode in slot $3$, the operands and the immediate in slots $4,\dots,10$, and $0$ everywhere else.
\item read the resulting $2^{\kbc}\times16$ table as $P$, the low $\kbc$ variables indexing the instruction and the high four its slot, bits low first.
\end{enumerate}
\begin{center}
\begin{tabular}{lllllllll}
\hline
Instruction & $3$ & $4$ & $5$ & $6$ & $7$ & $8$ & $9$ & $10$\\
\hline
\texttt{XOR64}         & $\opc{XOR64}$ & $o_A$ & $o_B$ & $o_C$ & $0$ & $0$ & $0$ & $0$\\
\texttt{MUL64}         & $\opc{MUL64}$ & $o_A$ & $o_B$ & $o_C$ & $0$ & $0$ & $0$ & $0$\\
\texttt{SET\_CONSTANT} & $\opc{SET}$ & $o$   & $k$ & $0$ & $0$ & $0$ & $0$ & $0$\\
\texttt{DEREF}         & $\opc{DRF}$ & $o_1$ & $o_2$ & $o_3$ & $f_{\pc}$ & $f_{\fp}$ & $0$ & $0$\\
\texttt{JUMP}          & $\opc{JMP}$ & $o_c$ & $o_d$ & $o_f$ & $0$ & $0$ & $0$ & $0$\\
\texttt{BLAKE2S}        & $\opc{B2S}$  & $o_{m_0}$ & $o_{m_1}$ & $o_{m_2}$ & $o_{m_3}$ & $o_{\mathit{cv}}$ & $o_{\mathit{out}}$ & $o_{\md}$\\
\texttt{ADD\_U64}       & $\opc{ADDU}$ & $o_A$ & $o_B$ & $o_C$ & $0$ & $0$ & $0$ & $0$\\
\texttt{MUL\_U64}       & $\opc{MULU}$ & $o_A$ & $o_B$ & $o_C$ & $0$ & $0$ & $0$ & $0$\\
\hline
\end{tabular}
\end{center}
So $P(z,1,1,0,0)$ is instruction $z$'s opcode, $P(z,0,0,1,0)$ its first operand, and $P(z,0,0,0,0)=P(z,0,0,1,1)=0$.

\subsection{Public input}\label{sec:e2e-pi}

The public input is the first four memory cells $\mem[\gen^{0}],\dots,\mem[\gen^{3}]$, four words both parties know, which hold the same values before the run and after it (\S\ref{sec:vm}). The verifier samples $r=(r_0,r_1)\in\E^{2}$ and, with nothing sent, pools the two claims
\[
  \mle{\meminit}(r_0,r_1,0,\dots,0)\;\qeq\;\sum_{b\in\cube{2}}\eq(r,b)\,\mem[\gen^{\,b_0+2b_1}]\;\qeq\;\mle{\memfin}(r_0,r_1,0,\dots,0),
\]
whose middle term, the multilinear extension of the four public words at $r$ ($r_0$ on the low bit), the verifier computes itself. By Fact~\ref{fact:eq} each side is the same sum over the four committed cells, so a mismatch leaves a nonzero multilinear polynomial in $r$, of total degree at most two, and each claim holds with probability at most $2/|\E|$ (Lemma~\ref{lem:sz}). The PCS opening binds the claims to the committed memories. The one on $\meminit$ is what makes the four words an input; the one on $\memfin$, an output. Every other cell of $\meminit$ is the prover's to choose (\S\ref{sec:vm}).

\subsection{Filling the tables}\label{sec:e2e-pad}

A table is proven over a power-of-two number of rows. The prover reaches one with \emph{padding rows}, whose clock is $0$. The bytecode carries, for each table and each of a few sizes, a block of that many instructions of the table's opcode followed by a \texttt{JUMP} back to the block's first instruction, which no program code jumps into. A padding row is a row at one of those instructions. Since $0\cdot\gen^{s}=0$, the states pushed and pulled around a block cancel among themselves, for any number of traversals: these are the closed walks of Proposition~\ref{prop:walk}, and by Corollary~\ref{cor:clock} they are the only ones. A traversal of the block of size $n$ adds $n$ rows to its table and one to \texttt{JUMP}'s, so every table can be brought to its next power of two exactly.

A padding row's memory accesses cost nothing either. The gap check \eqref{eq:gap} forces $\tsprev=0$ on them, so an access pulls $\tup{\dsep{MEM},a,0,v_{\mathrm{old}}}$ and pushes $\tup{\dsep{MEM},a,0,v_{\mathrm{new}}}$; the prover takes $v_{\mathrm{old}}=v_{\mathrm{new}}$ and the two cancel, whatever the address and the value. Its reads of the bytecode and of the range arrays are ordinary ones, with $\glo=\gen$ and $\ghi=1$, the two arrays' first entries. None of this bears on soundness: by Theorem~\ref{thm:rwmem} a tuple whose timestamp is zero meets no tuple of the run, so the padding rows constrain nothing and are constrained by nothing, which is what a padding row should be. The zero of $\K$, the one element that is not a power of $\gen$, plays the part of an \emph{is-real} selector, at the cost of no column and no constraint.

\subsection{Fiat Shamir instantiation}\label{sec:e2e-fiat-shamir}

TODO

\subsection{The unrolled protocol}\label{sec:e2e-unrolled}

Each reduction below adds evaluation claims on columns to a shared \emph{claim pool}, proved by the final PCS \textsc{Opening} phase.

\paragraph{Setup.}


The Fiat-Shamir initial state (\S\ref{sec:e2e-fiat-shamir}) is seeded by the public input, the bytecode and Flock's BLAKE2s R1CS matrices. The prover sends the memory log-size $\kmem$, the eight table log-heights $\tau_j$, the PCS rate and the final clock $\tsfinal\in\K$ (\S\ref{sec:state}); the verifier checks the sizes against the caps (\S\ref{sec:lookup}) and derives the layout of every stack (PCS and GKR leaves).

\paragraph{Commitment.}

\begin{enumerate}
\item \textbf{Prover.} Stack (\S\ref{sec:stacking}) every column, the two memories $\meminit$ and $\memfin$, the final timestamps $\tsfin$ (\S\ref{sec:memchan}), the three finalize counts (bytecode and the two range arrays, \S\ref{sec:rangecheck}), and Flock's three packed witnesses, $\qflock$ (\S\ref{sec:tab-blake2s}), $q_{\mathrm{add}}$ and $q_{\mathrm{mul}}$ (\S\ref{sec:tab-u64}), into $q$ and sends its commitment (Annex~\ref{annex:pcs}).
\item \textbf{Verifier.} Receives the commitment to $q$.
\end{enumerate}

\paragraph{Bus.}

\begin{enumerate}
\item \textbf{Verifier.} Samples and sends the fingerprint challenge $\alpha\in\E^{4}$ and the multiset challenge $\beta\in\E$ (\S\ref{sec:gp}).
\item \textbf{Prover.} Forms the push, pull, and \emph{count} leaf vectors (\S\ref{sec:leafstack}; the count leaves are the count columns themselves, \S\ref{sec:lookup}), stacked into blocks and padded with $1$s. Sends one bus root $R$, then the count root $R_c$.
\item \textbf{Prover \& Verifier.} Run the one batched GKR (\S\ref{sec:gkr}) over all three trees. For the bus, the verifier checks each side's product against $R$; one $R$ for both sides is the balance check (\S\ref{sec:gp}). For the counts it checks $R_c\neq0$ (no read self-cancels, \S\ref{sec:lookup}). The pass yields three leaf claims $\widetilde V^{s}_0(\zeta)$ at one shared $\zeta$.
\item \textbf{Prover \& Verifier.} Decompose each leaf claim (\S\ref{sec:leafstack}). The verifier forms every public part itself: the block selectors, the constant entries, the index and geometric columns (\S\ref{sec:idxcol}), and the program's whole share of the two bytecode blocks, one evaluation at $(\zeta,\alpha)$ (\S\ref{sec:e2e-bc}). Five blocks per side belong to no table: the state boundary, plus each array's (memory, bytecode, the two range arrays) seed (push side) or finalization (pull side). For their committed columns the prover sends one evaluation each, entering the claim pool. Every other block belongs to a table, and is discharged by the table sumcheck (\S\ref{sec:air}).
\end{enumerate}

\paragraph{Table sumcheck.} One run proves the eight tables' constraints together with their share of the bus (\S\ref{sec:air}).
\begin{enumerate}
\item \textbf{Verifier.} Samples and sends the batch challenge $\xi\in\E$. Its first $\nconstot=\sum_j\ncons{j}$ powers weight the constraints, a disjoint range per table; the last $\nside$ weight the bus forms, one per side and shared by every table. It draws no point: the batch runs at the bus point, $T_j$ tested at $\zeta_{<\tau_j}$.
\item \textbf{Prover \& Verifier.} Run the sumcheck, $\taumax=\max_j\tau_j$ rounds with $T_j$ joining at round $\taumax-\tau_j$, against the target the verifier derived from the three leaf claims. Its summand has individual degree three and the round check pins one of the four coefficients (\S\ref{sec:prelim}), so a round costs three $\E$-elements.
\item \textbf{Prover \& Verifier.} The prover sends one value per column of every table. The verifier rebuilds each $C_{j,i}$ and $B^{s}_j$ from them and checks that they reproduce the sumcheck's final value; the column values then enter the claim pool, $T_j$'s at $\fc_{<\tau_j}$. \texttt{BLAKE2S}'s eighteen value words are columns of its table but live in $\qflock$, so their claims route there (\S\ref{sec:tab-blake2s}); likewise the three value words of \texttt{ADD\_U64} and of \texttt{MUL\_U64}, in $q_{\mathrm{add}}$ and $q_{\mathrm{mul}}$ (\S\ref{sec:tab-u64}).
\end{enumerate}

\paragraph{Public input.}

\begin{enumerate}
\item \textbf{Verifier.} Sends $r=(r_0,r_1)\in\E^{2}$ and evaluates the multilinear extension of the four public words at $r$.
\item \textbf{Prover \& Verifier.} Pool that value as a claim on $\meminit$ and as a claim on $\memfin$, both at $(r_0,r_1,0,\dots,0)$; the prover sends nothing (\S\ref{sec:e2e-pi}).
\end{enumerate}

\paragraph{Flock validity.}

\begin{enumerate}
\item \textbf{Prover \& Verifier.} Run Flock's reduction (Annex~\ref{annex:flock}) three times, once per circuit and in this order: BLAKE2s over $\qflock$ (\S\ref{sec:tab-blake2s}), then the adder over $q_{\mathrm{add}}$ and the multiplier over $q_{\mathrm{mul}}$ (\S\ref{sec:tab-u64}). Each leaves one evaluation claim on its own packed witness.
\item \textbf{Prover \& Verifier.} The ring switching of \S\ref{sec:ringswitch}, with one map drawn for the three, reduces each bit-witness claim to a weighted-sum claim on its witness's region of the stack, with an MLE-friendly $\E$-valued weight (Definition~\ref{def:ipcs}); they join the pool.
\end{enumerate}

\paragraph{Opening.}

\begin{enumerate}
\item \textbf{Prover \& Verifier.} Each pooled claim is one evaluation of a column at a point, its value already supplied; via the stacking selectors (\S\ref{sec:stacking}) it becomes a weighted sum $\sum_w W_j(w)\,\widetilde q(w)=c_j$ over the stacked witness. (The three ring-switched claims arrive already in this form, each weight supported on its own witness's region.)
\item \textbf{Verifier.} Sends one batching challenge $\lambda\in\E$, drawn after every claim value is bound. Claim $j$ takes the weight $\lambda^{j-1}$, the ring-switched claims first and the pooled column claims after: the $J$ weights are distinct powers of the one challenge (Theorem~\ref{thm:rbr}).
\item \textbf{Prover \& Verifier.} Fold the $J$ claims into $W_\lambda=\sum_j\lambda^{j-1} W_j$ and $C_\lambda=\sum_j\lambda^{j-1} c_j$, and run the inner-product PCS opening on $\sum_w W_\lambda(w)\,\widetilde q(w)=C_\lambda$ (Annex~\ref{annex:pcs}), the verifier evaluating $W_\lambda$ itself. This one opening discharges every pooled claim; there is no separate reduction sumcheck.
\item \textbf{Verifier.} Accepts if every check above passed: the caps, the one bus root for both sides, the nonzero count root, every sumcheck's rounds and final value, flock's three reductions, and the PCS opening.
\end{enumerate}
